\documentclass[aspectratio=169]{beamer}
\usepackage[utf8]{inputenc}
\usepackage[italian]{babel}
\usepackage{graphicx}
\usepackage{amsmath}
\usepackage{amssymb}
\usepackage{listings}
\usepackage{xcolor}
\usepackage{tikz}
\usepackage{fontawesome5}
\usepackage{colortbl}
\usepackage{comment}
\usetikzlibrary{arrows.meta,shadows,shapes.geometric,positioning}

\usetheme{metropolis}
\usecolortheme{default}

\definecolor{primaryDark}{RGB}{45,55,72}
\definecolor{accentBlue}{RGB}{66,153,225}
\definecolor{darkGray}{RGB}{74,85,104}
\definecolor{lightGray}{RGB}{247,250,252}
\definecolor{mediumGray}{RGB}{160,174,192}
\definecolor{successGreen}{RGB}{72,187,120}
\definecolor{warningRed}{RGB}{245,101,101}

\setbeamercolor{palette primary}{bg=primaryDark,fg=white}
\setbeamercolor{palette secondary}{bg=accentBlue,fg=white}
\setbeamercolor{palette tertiary}{bg=darkGray,fg=white}
\setbeamercolor{frametitle}{bg=primaryDark,fg=white}
\setbeamercolor{title}{bg=primaryDark,fg=white}
\setbeamercolor{block title}{bg=primaryDark,fg=white}
\setbeamercolor{block body}{bg=lightGray,fg=darkGray}
\setbeamercolor{alerted text}{fg=warningRed}
\setbeamercolor{example text}{fg=successGreen}

\setbeamertemplate{frame numbering}[fraction]
\setbeamertemplate{footline}{
  \leavevmode%
  \hbox{%
  \begin{beamercolorbox}[wd=.5\paperwidth,ht=2.5ex,dp=1ex,left,leftskip=1em]{author in head/foot}%
    \usebeamerfont{author in head/foot}\textbf{Sandi Russo} \hspace{0.3cm} | \hspace{0.3cm} \textit{UniME - Scienze Informatiche}
  \end{beamercolorbox}%
  \begin{beamercolorbox}[wd=.5\paperwidth,ht=2.5ex,dp=1ex,right,rightskip=1em]{date in head/foot}%
    \usebeamerfont{date in head/foot}
    \insertframenumber{} / \inserttotalframenumber
  \end{beamercolorbox}}%
  \vskip0pt%
}

\setbeamertemplate{navigation symbols}{}

\lstset{
    basicstyle=\ttfamily\footnotesize,
    keywordstyle=\color{accentBlue}\bfseries,
    commentstyle=\color{successGreen}\itshape,
    stringstyle=\color{warningRed},
    breaklines=true,
    backgroundcolor=\color{lightGray},
    frame=none,
    numbers=left,
    numberstyle=\tiny\color{mediumGray},
    xleftmargin=1.5em,
    framexleftmargin=1em,
    rulecolor=\color{primaryDark},
    showstringspaces=false,
    tabsize=2
}

\tikzset{
    modernbox/.style={
        rectangle, rounded corners=2pt, 
        draw=primaryDark, thick, 
        fill=lightGray,
        drop shadow={shadow xshift=1.5pt, shadow yshift=-1.5pt, opacity=0.25}
    },
    technode/.style={
        circle, draw=primaryDark, thick,
        minimum size=1.1cm,
        drop shadow={shadow xshift=1pt, shadow yshift=-1pt, opacity=0.3}
    },
    modernArrow/.style={
        ->, thick, >=stealth, 
        color=primaryDark
    }
}

\title[Graph ML con Oracle]{Knowledge Graph e Graph Machine Learning}
\subtitle{Analisi Comparativa e Implementazione con Oracle}
\author{Sandi Russo}
\institute{Università degli Studi di Messina\\Corso di Laurea in Scienze Informatiche}
\date{}

\begin{document}

{
\setbeamertemplate{footline}{}
\begin{frame}
    \vspace{1cm}
    \begin{center}
        {\Large\textbf{Knowledge Graph e Graph Machine Learning}}
        
        \vspace{0.3cm}
        
        {\large Analisi Comparativa e Implementazione con Oracle}
        
        \vspace{1cm}
        
        \begin{tikzpicture}
            \node[draw=primaryDark, thick, rounded corners=3pt, fill=lightGray!30, 
                  text width=9cm, align=center, inner sep=0.5cm] {
                \small
                \begin{tabular}{cl}
                    \faUser & \textbf{Sandi Russo} \\[0.15cm]
                    \faUniversity & Università degli Studi di Messina \\[0.15cm]
                    \faGraduationCap & Scienze Informatiche \\[0.15cm]
                    \faCalendar & Anno Accademico 2024/2025 \\[0.15cm]
                    \faChalkboardTeacher & Relatore: docente Antonio Celesti \\
                \end{tabular}
            };
        \end{tikzpicture}
    \end{center}
\end{frame}
}

\begin{frame}{Agenda}
    \tableofcontents
\end{frame}

\section{Il Dataset: L'Ecosistema PharMeBINet}

\begin{frame}{L'Evoluzione dei Knowledge Graph Biomedici}
    \begin{columns}[T]
        \begin{column}{0.48\textwidth}
            \begin{alertblock}{\faTimesCircle \hspace{0.2cm} Il Punto di Partenza: Hetionet}
                \small
                \textbf{Hetionet (het.io)}
                \begin{itemize}
                    \item Storico Knowledge Graph biomedico.
                    \item Ottimo punto di partenza teorico, ma limitato.
                    \item Struttura obsoleta e instabile nell'importazione su DB moderni.
                \end{itemize}
            \end{alertblock}
        \end{column}
        
        \begin{column}{0.48\textwidth}
            \begin{block}{\faCheckCircle \hspace{0.2cm} L'Evoluzione: PharMeBINet}
                \small
                \textbf{The heterogeneous pharmacological network}
                \begin{itemize}
                    \item Fonde Hetionet con dati eterogenei moderni.
                    \item \textit{Data Engineering}: parsing, mapping di identificativi complessi e rimozione di sub-grafi ridondanti.
                \end{itemize}
            \end{block}
        \end{column}
    \end{columns}
    
    \vspace{0.3cm}
    \begin{exampleblock}{\faQuoteLeft \hspace{0.2cm} Riferimento Scientifico}
        \scriptsize
        Königs C., Friedrichs M., Dietrich T. \textit{The heterogeneous pharmacological medical biochemical network PharMeBINet}. Scientific Data. 2022;9(1): 393.
    \end{exampleblock}
\end{frame}

\begin{frame}{Anatomia di PharMeBINet}
    \begin{columns}[T]
        \begin{column}{0.5\textwidth}
            \textbf{\faProjectDiagram \hspace{0.2cm} Dati Generali del Grafo}
            \vspace{0.2cm}
            
            \begin{tikzpicture}[
                modernbox/.style={
                    rectangle, rounded corners=3pt, 
                    draw=primaryDark, thick, 
                    fill=lightGray,
                    drop shadow={shadow xshift=2pt, shadow yshift=-2pt, opacity=0.3}
                }
            ]
                \node[modernbox, fill=accentBlue!20, minimum width=4cm, minimum height=1cm] (nodes) at (0, 0) {
                    \textbf{Nodi}: 5.831.471
                };
                
                \node[modernbox, fill=successGreen!20, minimum width=4cm, minimum height=1cm] (edges) at (0, -1.5) {
                    \textbf{Relazioni}: 22.783.080
                };
                
                \node[modernbox, fill=warningRed!20, minimum width=4cm, minimum height=1cm] (labels) at (0, -3) {
                    \textbf{80 Labels} | \textbf{277 Edge Types}
                };
                
                \draw[thick, primaryDark, dashed] (nodes) -- (edges);
                \draw[thick, primaryDark, dashed] (edges) -- (labels);
            \end{tikzpicture}
        \end{column}
        
        \begin{column}{0.5\textwidth}
            \textbf{\faDatabase \hspace{0.2cm} Integrazione Dati Massiva}
            \small
            \vspace{0.2cm}
            
            PharMeBINet non è un singolo file, ma aggrega e armonizza oltre \textbf{40 database biomedici}, tra cui:
            \begin{itemize}
                \item \textbf{DrugBank} (Farmacologia)
                \item \textbf{ClinVar \& dbSNP} (Variazioni)
                \item \textbf{Gene Ontology} (Ontologie)
                \item \textbf{SIDER} (Effetti Collaterali)
                \item \textbf{DisGeNET} (Malattie)
            \end{itemize}
        \end{column}
    \end{columns}
\end{frame}

\begin{frame}{Il Valore Aggiunto: Complessità e Multi-Modalità}
    \begin{block}{\faDna \hspace{0.2cm} Molto più di un semplice grafo sociale}
        \small
        La maggior parte dei tutorial sul Graph ML utilizza dataset "giocattolo" e omogenei (es. reti di citazioni come Cora o Citeseer). PharMeBINet offre una rete \textbf{eterogenea e multi-modale}.
    \end{block}
    
    \vspace{0.3cm}
    \small
    \textbf{Perché questa complessità è essenziale per i test?}
    \begin{itemize}
        \item[\faNetworkWired] Permette di testare la capacità dei database di gestire query complesse (\textit{multi-hop}) tra entità di natura profondamente diversa (es. \textit{Gene $\rightarrow$ Proteina $\rightarrow$ Malattia $\rightarrow$ Farmaco}).
        \item[\faTags] Obbliga i modelli di Machine Learning a calcolare \textit{Node Embeddings} che tengano conto non solo della topologia pura, ma del \textit{tipo} di arco che collega i nodi.
    \end{itemize}
\end{frame}

\begin{frame}{Potenziale per il Graph Machine Learning}
    \small
    Le \textbf{Graph Neural Networks} applicate a questa struttura specifica sbloccano casi d'uso critici per la bioinformatica:
    
    \vspace{0.3cm}
    \begin{itemize}
        \item[\faMagic] \textbf{Link Prediction (Riposizionamento Farmaci)}: Se la GNN apprende il pattern topologico per cui un "Farmaco A" cura la "Malattia B", può prevedere archi futuri suggerendo cure inedite.
        \item[\faVial] \textbf{Previsione Effetti Collaterali}: Analizzando reti di interazione tra composti (\texttt{drug-drug interactions}), è possibile prevedere tossicità non documentate.
        \item[\faUsersCog] \textbf{Community Detection}: Raggruppare entità in modo non supervisionato per scoprire pattern biologici nascosti, come network di geni co-espressi che partecipano alla stessa via metabolica.
    \end{itemize}
\end{frame}

\begin{frame}{Un banco di prova "Real-World"}
    \begin{alertblock}{\faGlobe \hspace{0.2cm} La sfida dei dati reali}
        \small
        I dataset accademici standard sono perfettamente puliti. L'aggregazione di PharMeBINet porta con sé tutte le "sporcizie" tipiche degli scenari \textit{Enterprise}.
    \end{alertblock}
    
    \vspace{0.2cm}
    \small
    \textbf{Perché sceglierlo come Benchmark?}
    \begin{itemize}
        \item Mettere sotto stress i motori di importazione (Neo4j APOC vs Oracle \texttt{SQL*Loader}).
        \item Testare la tolleranza ai difetti dei parser (caratteri invisibili, \texttt{\textbackslash r}, EOF errati).
        \item Valutare come i diversi database (RDBMS vs Nativi) gestiscono le violazioni topologiche ("Dangling edges").
    \end{itemize}
\end{frame}

\section{Piani di Test: Analytics e Machine Learning}

\begin{frame}{Obiettivi di Testing: Graph Analytics}
    \begin{block}{\faChartBar \hspace{0.2cm} Analisi Topologica del Grafo}
        \small
        \begin{itemize}
            \item[\faBullseye] \textbf{Centrality (Centralità)}
            \begin{itemize}
                \scriptsize \item \textit{Algoritmi}: PageRank, Betweenness Centrality.
                \scriptsize \item \textit{Obiettivo}: Trovare i geni più influenti o i farmaci "hub".
            \end{itemize}
            
            \item[\faNetworkWired] \textbf{Community Detection}
            \begin{itemize}
                \scriptsize \item \textit{Algoritmi}: Louvain, Weakly Connected Components.
                \scriptsize \item \textit{Obiettivo}: Raggruppare entità biologiche con funzioni simili.
            \end{itemize}
            
            \item[\faRoute] \textbf{Path Finding}
            \begin{itemize}
                \scriptsize \item \textit{Algoritmi}: Shortest Path, A*.
                \scriptsize \item \textit{Obiettivo}: Trovare la catena più breve tra farmaco e sintomo.
            \end{itemize}
        \end{itemize}
    \end{block}
\end{frame}

\begin{frame}{Obiettivi di Testing: Graph Machine Learning}
    \begin{block}{\faRobot \hspace{0.2cm} Task di ML su Grafi}
        \small
        \begin{itemize}
            \item[\faVectorSquare] \textbf{Node Embeddings}
            \begin{itemize}
                \scriptsize \item Trasformare nodi in vettori preservando la struttura della rete.
            \end{itemize}
            
            \item[\faTags] \textbf{Node Classification}
            \begin{itemize}
                \scriptsize \item Classificare entità sconosciute basandosi sui loro "vicini".
            \end{itemize}
            
            \item[\faMagic] \textbf{Link Prediction}
            \begin{itemize}
                \scriptsize \item \textbf{Task critico}: Prevedere collegamenti mancanti (es. cure potenziali).
            \end{itemize}
            
            \item[\faObjectGroup] \textbf{Unsupervised Learning}
            \begin{itemize}
                \scriptsize \item Clustering dei vettori per scoprire nuove tassonomie biomediche.
            \end{itemize}
        \end{itemize}
    \end{block}
\end{frame}

\section{Implementazione Pratica: Setup e Ingestione Dati}

\begin{frame}[fragile]{Setup dell'Ambiente e Strumentazione}
    \begin{columns}[T]
        \begin{column}{0.48\textwidth}
            \begin{block}{\faLinux \hspace{0.2cm} Isolamento e Sviluppo}
                \small
                \begin{itemize}
                    \item OS: Ubuntu tramite \textbf{WSL}
                    \item Isolamento: \texttt{venv} Python
                    \item Credenziali sicure via \texttt{.env}
                \end{itemize}
            \end{block}
            
            \vspace{0.2cm}
            
            \begin{exampleblock}{\faCogs \hspace{0.2cm} Stack Librerie}
                \small
                \begin{itemize}
                    \item \textbf{oracledb}: Driver DB ufficiale
                    \item \textbf{oml4py \& pypgx}: ML \& Grafi Oracle
                    \item \textbf{neo4j}: Driver container
                \end{itemize}
            \end{exampleblock}
        \end{column}
        
        \begin{column}{0.48\textwidth}
            \textbf{\faTerminal \hspace{0.2cm} Inizializzazione Ambiente}
            \begin{lstlisting}[language=bash, basicstyle=\ttfamily\tiny]
# Creazione e attivazione venv
python3 -m venv .venv
source .venv/bin/activate

# Installazione dipendenze core
pip install oracledb neo4j python-dotenv pandas numpy jupyter oml4py pypgx
            \end{lstlisting}
        \end{column}
    \end{columns}
\end{frame}

\begin{frame}[fragile]{Inizializzazione di Neo4j}
    \textbf{\faDocker \hspace{0.2cm} Importazione del Dataset PharMeBINet}
    \small
    
    I dati originali pre-processati (file \texttt{.dump}) sono stati caricati in un container temporaneo per popolare l'istanza definitiva.
    
    \vspace{0.2cm}
    \begin{lstlisting}[language=bash, basicstyle=\ttfamily\tiny]
# 1. Assegnazione permessi per il mount dei volumi Docker
sudo chown -R $USER:$USER $HOME/neo4j

# 2. Importazione del dump tramite container temporaneo
docker run --rm \
  -v $HOME/neo4j/data:/data \
  -v $HOME/neo4j/import:/import \
  -e NEO4J_ACCEPT_LICENSE_AGREEMENT=yes \
  neo4j:5-enterprise \
  neo4j-admin database load neo4j --from-path=/import --overwrite-destination=true

# 3. Avvio del container ufficiale
docker start neo4j-gds
    \end{lstlisting}
\end{frame}

\begin{frame}[fragile]{L'Ostacolo della Versione Always Free}
    \begin{columns}[T]
        \begin{column}{0.48\textwidth}
            \begin{alertblock}{\faBan \hspace{0.2cm} Il Limite di Oracle}
                \small
                \begin{itemize}
                    \item Limite \textbf{hard-coded di 12 GB}.
                    \item TSV originali: oltre \textbf{30 GB}.
                    \item Impossibile caricare tutto il DB relazionale in Oracle.
                \end{itemize}
            \end{alertblock}
        \end{column}
        
        \begin{column}{0.48\textwidth}
            \begin{block}{\faLightbulb \hspace{0.2cm} Data Reduction Pipeline}
                \small
                Estrazione chirurgica (Bash/AWK) per estrarre un sub-grafo chiuso di circa 8 GB senza rischiare di avere nodi "orfani".
            \end{block}
        \end{column}
    \end{columns}
    
    \vspace{0.3cm}
    \textbf{\faTerminal \hspace{0.2cm} Pipeline di Riduzione su TSV}
    \begin{lstlisting}[language=bash, basicstyle=\ttfamily\tiny]
# 1. Isoliamo i primi 8 milioni di archi
head -n 8000001 edges.tsv > edges_medium.tsv

# 2. Estraiamo gli ID univoci e ritagliamo i nodi corrispondenti dal file sorgente
awk -F'\t' 'NR>1 {print $3; print $4}' edges_medium.tsv | sort -u > valid_nodes.txt
awk -F'\t' 'NR==FNR {v[$1]; next} FNR==1 || ($1 in v)' valid_nodes.txt nodes.tsv > nodes_medium.tsv
    \end{lstlisting}
\end{frame}

\begin{frame}[fragile]{Architettura e Importazione su Oracle DB}
    \begin{columns}[T]
        \begin{column}{0.48\textwidth}
            \begin{block}{\faDatabase \hspace{0.2cm} Preparazione Spazio}
                \small
                Allocazione di un \texttt{BIGFILE TABLESPACE} auto-espandibile per accomodare il traffico.
            \end{block}
        \end{column}
        
        \begin{column}{0.48\textwidth}
            \begin{alertblock}{\faExclamationTriangle \hspace{0.2cm} Scelta Ingegneristica}
                \small
                \textbf{Nessuna Foreign Key} in fase di creazione. Validare 8+ milioni di archi in inserimento avrebbe causato I/O insostenibile.
            \end{alertblock}
        \end{column}
    \end{columns}
    
    \vspace{0.4cm}
    \textbf{\faTractor \hspace{0.2cm} Travaso Alta Velocità (External Tables)}
    \begin{lstlisting}[language=SQL, basicstyle=\ttfamily\tiny]
-- Lettura diretta dai TSV su disco (tramite ORACLE_LOADER)
-- e inserimento massivo bypassando i log (Direct Path Load)
INSERT /*+ APPEND */ INTO pharme_edges (relationship_id, type, start_id, end_id) 
SELECT relationship_id, type, start_id, end_id 
FROM ext_pharme_edges;
COMMIT;
    \end{lstlisting}
\end{frame}

\begin{frame}{Il Limite dei 12GB e la Soluzione ML-Oriented}
    \begin{alertblock}{\faBan \hspace{0.2cm} Il Crash per i Campi Testuali}
        \small
        Il campo \texttt{properties} (CLOB) conteneva JSON testuali massicci, decuplicando lo spazio di archiviazione per via dell'overhead di Oracle e facendo scattare il blocco dei 12GB nonostante la riduzione a 8 milioni di archi.
    \end{alertblock}
    
    \vspace{0.2cm}
    
    \begin{exampleblock}{\faBrain \hspace{0.2cm} La Soluzione orientata al Machine Learning}
        \small
        \begin{itemize}
            \item Le \textbf{Graph Neural Networks} apprendono prevalentemente dalla \textbf{Topologia} e dalle \textbf{Labels}, non dal testo libero.
            \item \textbf{Azione:} Eliminazione dei CLOB testuali in fase di importazione, mantenendo solo ID strutturali, Label semantiche e Target relazionali.
        \end{itemize}
    \end{exampleblock}
    
    \vspace{0.2cm}
    \centering
    \faCheckCircle \hspace{0.1cm} \textbf{Risultato del Caricamento:} \textbf{5.831.471} Nodi e \textbf{22.783.080} Archi integrati.
\end{frame}

\begin{frame}[fragile]{Creazione del Property Graph Logico}
    \textbf{\faProjectDiagram \hspace{0.2cm} Da Tabelle Relazionali a Grafo (SQL/PGQ)}
    \small
    
    I vincoli logici e topologici, omessi fisicamente nelle tabelle relazionali per performance, vengono imposti rigorosamente nel costrutto del Grafo Logico su Oracle PGX.
    
    \vspace{0.2cm}
    \begin{lstlisting}[language=SQL, basicstyle=\ttfamily\scriptsize]
CREATE PROPERTY GRAPH pharme_graph
  VERTEX TABLES (
    pharme_nodes KEY (node_id)
      PROPERTIES (node_id, labels, name, identifier)
  )
  EDGE TABLES (
    pharme_edges KEY (relationship_id)
      SOURCE KEY (start_id) REFERENCES pharme_nodes (node_id)
      DESTINATION KEY (end_id) REFERENCES pharme_nodes (node_id)
      PROPERTIES (relationship_id, type)
  );
    \end{lstlisting}
\end{frame}

\section{Zero-Difference Benchmark e Data Forensics}

\begin{frame}[fragile]{Allineamento Dati (Zero-Difference Benchmark)}
    \textbf{\faBalanceScale \hspace{0.2cm} Garantire la Validità Scientifica}
    \small
    
    Per un benchmark ML equo, Oracle e Neo4j devono operare sugli \textbf{stessi identici byte}. I dati snelliti di Oracle sono stati esportati in file CSV puri (\texttt{sqlplus SET MARKUP CSV}) per sovrascrivere l'istanza originale di Neo4j.
    
    \vspace{0.2cm}
    \begin{block}{\faBug \hspace{0.2cm} Risoluzione Problemi di Importazione (Neo4j APOC)}
        \small
        \begin{itemize}
            \item \textbf{Pulizia Header:} Header CSV sporco rimosso via bash (\texttt{sed}).
            \item \textbf{Deadlock Multi-thread:} L'importazione massiva di archi con \texttt{parallel: true} causava eccezioni \textit{ExclusiveLock} sui nodi centrali (hub).
            \item \textbf{Fix Applicato:} Fallback all'esecuzione thread-safe (\texttt{parallel: false}) sacrificando la velocità per garantire l'integrità transazionale.
        \end{itemize}
    \end{block}
\end{frame}

\begin{frame}{Data Forensics: Indagine sulle Discrepanze}
    \begin{columns}[T]
        \begin{column}{0.48\textwidth}
            \begin{block}{\faGhost \hspace{0.2cm} Il Nodo Fantasma (+1)}
                \small
                \textbf{Problema}: Neo4j riportava 1 nodo in più di Oracle.\\
                \textbf{Causa}: Il parser APOC ha interpretato l'EOF come riga, generando un nodo generico \texttt{BaseNode} vuoto.\\
                \textbf{Fix}: 
                \texttt{\tiny MATCH (n:BaseNode) WHERE size(labels(n)) = 1 DETACH DELETE n;}
            \end{block}
        \end{column}
        
        \begin{column}{0.48\textwidth}
            \begin{alertblock}{\faUnlink \hspace{0.2cm} Gli Archi Orfani (-8.592)}
                \small
                \textbf{Problema}: Neo4j ha scartato 8.592 archi effettivamente presenti nel CSV.\\
                \textbf{Causa}: Presenza di Carriage Return (\texttt{\textbackslash r}) negli ID esportati, che impedivano il pattern-matching topologico di Neo4j.
            \end{alertblock}
        \end{column}
    \end{columns}
\end{frame}

\begin{frame}{Filosofia RDBMS vs Native Graph: La Prova del Nove}
    \begin{block}{\faCogs \hspace{0.2cm} Comportamento Fisico vs Logico}
        \small
        \begin{itemize}
            \item \textbf{Oracle DB (Fisico)}: Salva le righe tollerando inconsistenze (\texttt{\textbackslash r}) e "dangling edges", agendo come pura tabella relazionale.
            \item \textbf{Neo4j (Nativo)}: Pulisce i caratteri e scarta gli archi orfani, non ammettendo violazioni strutturali della rete.
        \end{itemize}
    \end{block}
    
    \vspace{0.2cm}
    
    \begin{exampleblock}{\faCheckDouble \hspace{0.2cm} L'attivazione del Motore Graph (PGX)}
        \small
        Interrogando il Grafo Logico (\texttt{pharme\_graph}) tramite SQL/PGQ, Oracle \textbf{filtra il rumore}:
        \vspace{0.1cm}
        \centering
        \texttt{\scriptsize SELECT COUNT(*) FROM GRAPH\_TABLE ( pharme\_graph MATCH ()-[r]->() ... )}
        \vspace{0.1cm}
        
        \textbf{Risultato per entrambi}: \textbf{22.783.080} archi validi. \\
        \textit{Il motore PGX di Oracle applica lo stesso rigoroso filtro topologico del motore nativo di Neo4j.}
    \end{exampleblock}
\end{frame}

\section{Bilancio del Progetto e Roadmap Futura}

\begin{frame}{Considerazioni Finali: Il Data Engineering}
    \begin{block}{\faCheckSquare \hspace{0.2cm} Obiettivo Infrastrutturale Raggiunto}
        \small
        Nonostante i pesanti limiti della versione \textit{Free Tier}, è stato allestito con successo un ambiente di testing \textbf{Zero-Difference} su un dataset biomedico massivo (\textasciitilde 23M relazioni).
    \end{block}
    
    \vspace{0.3cm}
    \small
    \textbf{Lezioni Apprese:}
    \begin{itemize}
        \item[\faTractor] L'inserimento massivo RDBMS richiede il bypass temporaneo dei vincoli (\textit{No FK, Direct Path Load}).
        \item[\faBrain] Nel Graph ML, la \textbf{topologia} conta più del testo descrittivo: eliminare JSON pesanti ha salvato il progetto senza compromettere i task predittivi.
        \item[\faSearch] Il motore logico di Oracle (PGX) converte perfettamente l'anarchia relazionale nel rigore topologico nativo tipico di Neo4j.
    \end{itemize}
\end{frame}

\begin{frame}{Considerazioni Finali: Oracle vs Neo4j}
    \begin{columns}[T]
        \begin{column}{0.48\textwidth}
            \begin{alertblock}{\faBalanceScale \hspace{0.2cm} Il Trade-off di Oracle}
                \small
                L'approccio di Oracle richiede una \textbf{curva di apprendimento molto più ripida} rispetto a Neo4j. Mantenere l'allineamento tra tabelle relazionali e astrazione logica (PGX) raddoppia la complessità del setup iniziale.
            \end{alertblock}
        \end{column}
        
        \begin{column}{0.48\textwidth}
            \begin{exampleblock}{\faRocket \hspace{0.2cm} Il Vantaggio Strategico}
                \small
                Il compromesso è giustificato: la versione \textit{Always Free} sblocca algoritmi \textbf{Enterprise-grade} e librerie native di ML (\texttt{OML4Py}) che in ecosistemi concorrenti richiederebbero licenze commerciali costose.
            \end{exampleblock}
        \end{column}
    \end{columns}
\end{frame}

\begin{frame}{Prospettive: Benchmark e Algoritmi}
    \small
    Le fondamenta ingegneristiche gettate in questa fase permettono di passare all'esecuzione del \textbf{Graph Machine Learning} puro:
    
    \vspace{0.3cm}
    \begin{itemize}
        \item[\faChartLine] \textbf{Benchmark Computazionale}: Esecuzione di \textit{PageRank} e \textit{Louvain} su 22 milioni di archi per confrontare i tempi di inferenza e il consumo RAM tra il motore PGX in-memory e Neo4j GDS.
        \item[\faProjectDiagram] \textbf{Graph Neural Networks Avanzate}: Applicazione dei modelli \textit{GraphSAGE} e \textit{Graph Attention Networks (GAT)} tramite l'interfaccia Python \texttt{PyPGX}.
        \item[\faDna] \textbf{Link Prediction Biomedica}: Costruire un task di test reale per prevedere probabili iterazioni farmaco-malattia non presenti nel dataset.
    \end{itemize}
\end{frame}

\begin{frame}{Lavori Futuri: Task Predittivi con KG Embeddings}
  \begin{block}{\faTags \hspace{0.2cm} Node Classification tramite Embeddings}
    \small
    Utilizzo di tecniche di generazione di Knowledge Graph (KG) Embeddings per mappare i nodi in uno spazio vettoriale a bassa dimensionalità.
    \begin{itemize}
      \item \textbf{Caso d'Uso (PharMeBINet)}: Prevedere la categoria di appartenenza di un nuovo composto o inferire la funzione biologica di un gene non annotato, basandosi sulla sua posizione nello spazio latente.
    \end{itemize}
  \end{block}

  \vspace{0.1cm}
  \begin{exampleblock}{\faLink \hspace{0.2cm} Link Prediction basata su Relazioni}
    \small
    Generazione di embeddings specifici non solo per i nodi, ma anche per le \textbf{relazioni} (archi), al fine di calcolare la probabilità di esistenza di collegamenti mancanti.
    \begin{itemize}
      \item \textbf{Caso d'Uso (PharMeBINet)}: Riposizionamento dei farmaci (\textit{Drug Repurposing}) prevedendo nuovi archi \texttt{TREATS} tra farmaci esistenti e malattie, o scoperta di interazioni avverse inedite.
    \end{itemize}
  \end{exampleblock}
\end{frame}

\begin{frame}{Lavori Futuri: Tecniche di Embedding Relazionale}
  \small
  Per implementare i task di Link Prediction e Node Classification, si prevede di testare e confrontare i seguenti modelli di \textbf{Knowledge Graph Embedding}:

  \vspace{0.3cm}
  \begin{itemize}
    \item[\faArrowRight] \textbf{TransE (Translational Embeddings)}
    \begin{itemize}
      \scriptsize \item Interpreta le relazioni come \textit{traslazioni} nello spazio vettoriale ($h+r\approx t$).
      \scriptsize \item \textbf{Vantaggio}: Intuitivo ed efficiente, ideale per modellare gerarchie (\textit{is-a}) o pattern 1-a-1.
    \end{itemize}

    \vspace{0.2cm}
    \item[\faExpand] \textbf{DistMult}
    \begin{itemize}
      \scriptsize \item Modello bilineare basato sulla fattorizzazione tensoriale con matrici diagonali.
      \scriptsize \item \textbf{Vantaggio}: Cattura in modo eccellente le interazioni \textit{simmetriche} (es. interazioni proteina-proteina).
    \end{itemize}

    \vspace{0.2cm}
    \item[\faSuperscript] \textbf{ComplEx (Complex Embeddings)}
    \begin{itemize}
      \scriptsize \item Estende DistMult operando nel dominio dei \textbf{numeri complessi}.
      \scriptsize \item \textbf{Vantaggio}: Modella efficacemente le relazioni \textit{asimmetriche} (es. un gene che regola un altro gene, ma non viceversa), dinamica fondamentale in biologia.
    \end{itemize}
  \end{itemize}
\end{frame}

\begin{frame}{Risorse e Riferimenti}
    \footnotesize
    \setlength{\itemsep}{2pt}
    \setlength{\parskip}{0pt}
    
    \textbf{\faBook \hspace{0.2cm} Oracle Database \& Graph}:
    \begin{itemize}
        \setlength{\itemsep}{1pt}
        \item Property Graph Docs: \texttt{docs.oracle.com/.../property-graph/}
        \item OML4Py: \texttt{docs.oracle.com/.../oml4py/}
        \item Graph Studio: \texttt{docs.oracle.com/.../graph-studio...}
    \end{itemize}
    
    \textbf{\faDatabase \hspace{0.2cm} Ecosistema \& Alternative}:
    \begin{itemize}
        \setlength{\itemsep}{1pt}
        \item PharMeBINet Dataset: \texttt{pharmebi.net}
        \item Neo4j: \texttt{graphacademy.neo4j.com}
        \item ArangoDB: \texttt{arangodb.com/docs/}
    \end{itemize}
    
    \textbf{\faPython \hspace{0.2cm} Librerie Machine Learning}:
    \begin{itemize}
        \setlength{\itemsep}{1pt}
        \item PyTorch Geometric: \texttt{pytorch-geometric.readthedocs.io}
        \item DGL: \texttt{dgl.ai}
    \end{itemize}
\end{frame}

{
\setbeamertemplate{footline}{}
\begin{frame}[standout]
    \Huge Grazie per l'attenzione!
\end{frame}
}

\end{document}